\section{PA194LC\_AX11H and the requested consistency equivalence}
\label{sec:lc-ax11h}

The final construction applies the same loop-continuation transformation
to the revised AX11H packet's primary machine. Denote that parent by
$M_H$ and the new machine by $\MLH$. The parent dense table is bound by
\begin{quote}\small
\path{c5dfe1026d077d3e062eca11edb789597bfbbced4f6040f8828ae5d467b1167a},
\end{quote}
and the new dense table by
\begin{quote}\small
\path{5f10f071e7118ff33aba822e05bae0331b4c78508656c3e945c31b0c6788f9c5}.
\end{quote}
The full new table is printed in Appendix~\ref{app:lc-ax11h-table} and
provided in \path{additional/PA194LC_AX11H/machine.tm}.

\subsection{Exact construction and completed checks}

The parent uses the four-field AX11 command, an eleven-entry saved prefix,
visible no-push MP, and the H halt test: Shift halts precisely when the
cached theorem is code one. The loop transformation leaves that command
protocol and its numerical calculus unchanged. All 3,671 resolved
instructions, including every register operand and successor, are
identical. The 61 recurring constructor loops change association and
replace their depth-two continuation by a depth-three continuation;
305 primitive addresses change.

Both the raw and final folded native tables differ in exactly the three
rows displayed in Section~\ref{sec:construction}: $(139,0)$ becomes
$0R99$, $(139,1)$ becomes $1R190$, and $(190,1)$ becomes $1L73$.
Here the parent of the patch is $M_H$, not the original primary $M_0$.
All seven contextual folds remain valid, giving $134+67-7=194$ working
controls and 388 transition rows. Initial and terminal controls are 0
and 194 in the dense numbering. The complete blank startup takes 57,922
transitions in both implementations.

\begin{center}
\begin{tabular}{@{}lrr@{}}
\toprule
Independent check & Parent AX11H & New LC\_AX11H\\
\midrule
Resolved source instructions & 3,671 & 3,671\\
Universal branch outcomes & 4,778 & 4,778\\
Symbolic operations & 1,101,836 & 1,101,043\\
Maximum operations in one case & 403 & 403\\
Minimum duration constant & 3 & 3\\
List lemmas derived from literal rows & 4 & 4\\
Whole-tape startup transitions & 57,922 & 57,922\\
\bottomrule
\end{tabular}
\end{center}

These are fresh full literal-row replays of both tables. The separate
construction checker also passed all 4,778 outcomes, using 1,514,306
symbolic operations, maximum 506. Both independent address resolvers
agree; complete source identity, native/dense conjugacy, simultaneous
fold cases, and full startup tape equality pass. The changed successful
loop iteration costs exactly six additional physical transitions at each
of its 61 occurrences. All five newly delivered LC graphs are pairwise
distinct and match none of the 110 prior graphs in the latest audited
snapshot. GAX11's separate duplicate is deliberately excluded.

\subsection{An explicit relative iPRA theorem}

Let $H_M(t)$ again mean halted by time $t$ in absorbing terminal semantics.
The finite coefficient audit of all 4,778 contracts in each table gives
the same envelope as in Section~\ref{sec:quadratic-transfer}:
\[
 3\le k\le534+8g+4\sum_{i<9}r_i+2a_0+2\operatorname{width}(A).
\]
The actual maximum vector is
$(534,8,4,4,4,4,4,4,2,4,4,2,2)$.
The exact startup frames have $g=0$, logical registers $(2,0,\ldots,0)$,
first spare $a_0=1$ and 25 further unit spares. Consequently the same
primitive updates give the same bound $k(s_j)\le694+26j$.

\begin{theorem}[Relative AX11H simulation]
With
\[
 \boxed{\Theta(t)=57\,922+13(t+1)^2+681(t+1)
                    =13t^2+707t+58\,616,}
\]
the following implications are provable in iPRA:
\[
 \forall t\,[H_{M_H}(t)\to H_{\MLH}(\Theta(t))],\qquad
 \forall t\,[H_{\MLH}(t)\to H_{M_H}(\Theta(t))].
\]
In particular iPRA proves equivalence of the two blank nonhalting
sentences.
\end{theorem}

\begin{proof}
Use the canonical zipper, explicit augmented source iteration, common
absorbing logical halt extension, Boolean prefix $B(n)$ and exact
clock $T(n)$ from Section~\ref{sec:quadratic-transfer}, now instantiated
with the 3,671-instruction AX11H source. Each local contract is proved by
the same literal-transition, unary-run and finite-list rules. The
independent checks establish all source targets and exact frame updates
for these two specific tables, so the quantifier-free synchronization
invariant applies with startup constant 57,922.

A physical halt by $t$ gives, by the bounded first-hit argument, a
common source halt-site index $n\le t$. In the other implementation,
the first $n$ normal instructions and its final halt instruction take
at most
\[
57\,922+\sum_{j=0}^{n}(694+26j)
=57\,922+13(n+1)^2+681(n+1).
\]
Monotonicity and absorbing halting yield the displayed witness at
$\Theta(t)$. All constructions and bounded Boolean invariants are
primitive recursive. No unbounded witness search or Markov principle
is used. The argument in both directions is identical after exchanging
the tables. A self-contained expanded version is included in the new
machine's standalone \path{proof.tex} and PDF.
\end{proof}

\subsection{Importing the revised absolute calibration}

The original upstream archive at the final snapshot is preserved without
modification as
\path{reference/ax11h_calibration_603c58c.zip}; its SHA-256 is
\begin{quote}\small
\path{93ee082ef035b2d32e71d6af5239b25b4bea1773f03bea4d2c39810a7fe15075}.
\end{quote}
Its \path{proof/consistency_equivalence.tex} states and proves, using its
separate completion papers and the pinned generic transfer results,
\[
 \ipra\vdash\neg\Halt(M_H)\leftrightarrow\Con(\pa).
 \tag{AX11H calibration}
\]
Here $\Con(\pa)$ is precisely the packet's
\texttt{nativePAConsistency}. This negates \texttt{nativePARefutation}
for Foundation dependency
\path{32e1a0956a8622fad067328ca1959729a7634428}.
Foundation falsity code 13 and numerical-calculus contradiction code 1
are connected by explicit proof translations; they are not identified.
\cite{latest-ax11h}

The imported argument supplies more than an assertion that a standard
simulation can be internalized. Its Foundation adapter gives explicit
proof opening/closing constructors, bounded induction-axiom witnesses,
and the correction for the changed finite order axiom. Its term and
formula compilers use acyclic primitive-recursive supplier strata and
bounded demand trees, with Boolean correctness invariants for the
actual requests, instead of full-PA induction over all environments.
Its source-command paper binds the literal arithmetic templates and
four-field parser, and its physical paper proves frame closure and
the symbolic rules uniformly. The final paper composes the two
all-model $I\Sigma_1$ directions with the existing one-way transfer
endpoints. The packet does not implement the entire large PA proof
compiler as executable code and does not supply a new end-to-end
Lean or Rocq build.

The focused review retained in \path{verification/ax11h_absolute_audit.txt}
found no identified gap in those completion arguments. The upstream
quick verifier passed all 1,186 manifest files, byte-identical
reconstruction of both parent designs, source-command and physical-frame
audits, fixed proof constructors, 119 current-Foundation coding
comparisons, and the premise-exchange certificate. That quick mode
explicitly does not verify the entire arithmetic derivation; the full
literal physical replays for $M_H$ and $\MLH$ were performed separately
as reported above.

\begin{theorem}[The 194-state fallback, importing the AX11H theorem]
Using the revised packet's written AX11H calibration as an imported
mathematical theorem,
\[
 \boxed{\ipra\vdash\neg\Halt(\MLH)\leftrightarrow\Con(\pa).}
\]
\end{theorem}
\begin{proof}
The relative theorem gives
$\neg\Halt(\MLH)\leftrightarrow\neg\Halt(M_H)$ in iPRA.
Compose it intuitionistically with the imported AX11H calibration.
\end{proof}

This explicitly identifies the inherited mathematical dependency.
The new contribution is the different literal machine and its exact,
positive-duration iPRA transfer, with complete source and tape evidence.
The corollary is a written proof using that supplied calibration, not
a claim that the unfinished primary-machine Lean development has acquired
a new kernel-checked endpoint.
